% 06e-actuator.tex 中文对照版（译自 sections/06e-actuator.tex；不加 \label）
\subsection{执行器闸门：面向工具调用的 STPA（System-Theoretic Process Analysis，系统理论过程分析）}

第 4 节的闸门管辖的是模型\emph{读入}什么。工具调用是另一个问题：它们会实际执行并产生副作用。一次被软化的拒绝只损失一轮对话；一次被删除的目录则损失数据。工具副作用不可逆，且其危害并不在注入侧——模型可能在一个冲突尚未解决时、在上下文用量处于前馈阈值时、或在经历了一轮社会工程之后，发出一次调用。因此我们把同样的“因果介入”（causal interposition）思想扩展到“工具请求”与“其执行”之间的区间，并使用一个独立的不安全控制动作（unsafe control action, UCA）命名空间——\emph{T-UCA}——它既不修改也不合并注入侧的 UCA-01--08。两道闸门是并列的：注入侧前置审计，然后工具闸门，然后执行。

\paragraph{范围，作为一种越权声明。} 当前项目策略冻结了工具闭环方向的工作。本扩展是一次\emph{显式越权解冻}（explicit override），而不是证据驱动的解冻：事先声明的触发条件——真实日志显示至少十个任务因外部调用不可用而失败，加上已证明无法回滚工具错误——\emph{并未}满足。我们报告它，是因为它已实现且可测试，但它的存在并不构成“触发条件已满足”的证据。三项内容仍保持冻结：多步工具规划、事务性回滚或补偿，以及把工具结果写入长期记忆。

\paragraph{T-UCA 矩阵。} 十六条规则——契约版本 \texttt{t-uca-1.1}——分为六组。\emph{状态耦合}（四条）：在多样性降维或未解决冲突下的写调用、临界相位下的写、情感基线饱和下的不可逆写入、以及 shift 相位下的破坏性调用。\emph{数据保护}（一条）：明文凭证。\emph{变更安全}（四条）：没有精确确认的破坏性删除或覆盖、写入冻结或受保护路径、冻结区校验失败后的写调用、以及路径逃逸。\emph{网络与执行}（两条）：非白名单外发与任意代码执行。\emph{形状与注册}（四条）：超大或递归 payload、schema 或分隔符违规、未登记工具、未声明能力。\emph{可审计性}（一条）：审计上下文不完整。十三条规则判定为 block，三条为 degrade：degrade 不执行原调用，而是要求通过宿主审批 seam（approval seam）获得显式人工确认，在没有审批通道时回落为 block。没有任何规则返回 \emph{allow}；allow 仅对“未命中任何规则”的只读调用作为默认，未知或不可判定的调用默认 block。

\paragraph{为什么状态重要。} 与只读取 prompt 的注入侧闸门不同，工具闸门消费宿主状态：多样性降维、未解决冲突、前馈监控的相位、积分饱和、冻结层校验、已声明能力，以及审计上下文。同一个调用在不同状态下不是同一个行为：在冲突未解决时删除文件，与澄清之后再删除，并不相同；而临界相位下的写入，正是不可逆错误最难归属的地方。

\paragraph{三个接入点。} 我们明确说明哪些接线是前置阻断式的。（i）在引擎的参考循环中，闸门位于解析之后、执行之前，因此被 block 的调用永远不会执行。（ii）在宿主插件中，\texttt{guardTool(spec, stateProvider)} 包装已登记的工具，使经由包装器路由的调用先被判定：这是前置阻断，但只对经由它路由的工具有效。（iii）宿主暴露 \texttt{tool/call} 与 \texttt{tool/result} 会话事件，但没有 pre-dispatch 钩子，因此插件的默认路径是\emph{事后检测}（detective）：它判定调用、记录一条 \texttt{tool\_gate} 审计事件（工具、参数摘要、判定、\texttt{ucas} 中命中的标识符、策略版本、状态快照），并在下一轮注入一条强澄清提示。检测不是阻断，我们如此标注。

\paragraph{证据。} Python 套件覆盖二十个构造的危险场景（全部被阻断）、八个良性场景（无一被阻断）、状态门控与 fail-closed 用例，以及一个二十例的 Python--JavaScript 检查，其中判定、命中标识符与降级字段全部一致。JavaScript 模拟运行六十五项检查；Python 套件报告 79 条通过的断言。参考演示输出本文引用的 JSON（\texttt{paper/ablation/results/actuator-demo-20261027.json}）：十二个脚本化调用，四个被允许并执行，八个被阻断且无一执行，命中标识符中涉及九条不同规则。它使用一个 mock 执行器，零模型调用、零真实写入，因此这些是展示“先判定、后执行”的场景计数，而不是拦截率。同一次运行还评估了一个盲写的十项留出集——其作者未读过规则表、实现或测试——冻结于 \texttt{paper/ablation/tool/heldout-tool-calls-v1.jsonl}（SHA-256 \texttt{f68fc1b3f19fc37e43b9856096d29b9a55de7dee9bf664a9b9a9592e070576d4}）；仅给参数时十次调用全部被阻断，加上上下文时同样十次全部被阻断，且其端点均为占位域名，因此命中不可能来自域名巧合。我们报告判定，而不报告正确率：每个调用\emph{是否应当}被阻断，仍待裁定。网关是增量的：注入侧规则未被触碰，且三项早于本工作的套件失败与环境相关、没有变化。

\paragraph{这些数字不是什么。} 我们不把这些计数换算成防御率：构造场景由设计者选定，而留出集只有一位盲写作者的十个案例。闸门从不执行 payload，因此其攻击面是解析，其误杀率尚未标定；恢复路径是一条强提示加人工确认，而不是静默拒绝。

\paragraph{边界。} 第一，没有宿主 pre-dispatch 钩子，插件无法阻止任意的宿主工具；真正的宿主级前置阻断是一项上游请求，不是已交付能力。第二，如果工具在网关之外执行——另一个插件，或绕过包装器的宿主路径——闸门就是空转的：它是一个可以被接入的控制点，而不是全系统强制。第三，上文列出的被排除项仍然冻结，因此本文不对闭环工具行为作任何主张。
